% !TeX program = lualatex
% Compile twice: lualatex open_problems_500.tex
% All references and prose are embedded; no BibTeX database or external inputs.
\documentclass[11pt,a4paper]{article}
\usepackage[margin=24mm,headheight=14pt]{geometry}
\usepackage{fontspec}
\setmainfont{Latin Modern Roman}
\setsansfont{Latin Modern Sans}
\setmonofont{Latin Modern Mono}
\usepackage{amsmath,amssymb,mathtools}
\usepackage{unicode-math}
\setmathfont{Latin Modern Math}
\usepackage{microtype}
\usepackage{enumitem,xurl,needspace}
\usepackage[dvipsnames]{xcolor}
\usepackage{hyperref}
\hypersetup{unicode=true,colorlinks=true,linkcolor=MidnightBlue,urlcolor=MidnightBlue,
 pdftitle={Research notebook: Problems 40--49},pdfauthor={Source-annotated compilation},bookmarksnumbered=true}
\usepackage{bookmark}
\usepackage{tocloft}
\setlength{\cftsecnumwidth}{3em}
\cftsetpnumwidth{2.5em}
\cftsetrmarg{4em}
\usepackage{titlesec}
\titleformat{\section}{\Large\bfseries\raggedright}{\thesection}{0.7em}{}
\usepackage{fancyhdr}
\pagestyle{fancy}\fancyhf{}
\fancyhead[L]{\small\sffamily Mathematical problem statements}
\fancyhead[R]{\small Source edition 22}
\fancyfoot[C]{\thepage}
\setlength{\parindent}{0pt}
\setlength{\parskip}{5pt plus 1pt}
\setlength{\emergencystretch}{3em}
\setcounter{tocdepth}{1}
\setcounter{secnumdepth}{1}
\setlist[itemize]{leftmargin=3em,itemsep=5pt,topsep=4pt,parsep=0pt}
\allowdisplaybreaks
\newcommand{\N}{\mathbb{N}}
\newcommand{\Z}{\mathbb{Z}}
\newcommand{\Q}{\mathbb{Q}}
\newcommand{\R}{\mathbb{R}}
\newcommand{\C}{\mathbb{C}}
\newcommand{\F}{\mathbb{F}}
\newcommand{\A}{\mathbb{A}}
\newcommand{\PP}{\mathbb P}
\newcommand{\E}{\mathbb{E}}
\newcommand{\eps}{\varepsilon}
\newcommand{\dd}{\,\mathrm d}
\newcommand{\sref}[2]{\hyperlink{source:#1:#2}{[S#2]}}
\newcommand{\eref}[2]{\hyperlink{extra:#1:#2}{[E#2]}}
% Turn the next switch off for a shorter reading copy. The quotations preserve
% every exactTarget field, including qualifications omitted from a short summary.
\newif\ifshowcatalogue
\showcataloguetrue
\newcommand{\cataloguescope}[1]{\ifshowcatalogue
 \paragraph{Exact scope in the supplied catalogue.}
 \begin{quote}\small #1\end{quote}\fi}

% Research additions dated 27 September 2026; compile twice with LuaLaTeX.
\numberwithin{equation}{section}
\begin{document}
\setcounter{section}{39}
% ================================================================
% problemId: problem.hilbert-sixteenth-problem-second-part
% source releaseRank: 40; source releaseStatus: open
% exactTarget SHA-256: 81dcde61465fa263d9218a450deb18aa55952f304b4113c688911a39b3528efe
\clearpage
\section{\texorpdfstring{Hilbert's Sixteenth Problem, Second Part}{Hilbert's Sixteenth Problem, Second Part}}\label{40}
\subsection{Definitions and mathematical statement}
A planar polynomial vector field of degree at most $n$ is
\[
 \dot x=P(x,y),\qquad\dot y=Q(x,y),\qquad P,Q\in\R[x,y],\quad\max(\deg P,\deg Q)\le n.
\]
A limit cycle is an isolated periodic orbit, not an arbitrary member of a continuous family of periodic trajectories. Let $H(n)$ be the supremum of the number of limit cycles over this class. The problem asks for finiteness of $H(n)$, its maximal value, and the possible relative configurations of these cycles for each degree. A bound depending on the individual coefficients is not a uniform degree bound.
\par\smallskip\textit{Formulation and concept references:} \sref{40}{1}.
\cataloguescope{For every degree n >= 2, prove finiteness and determine the maximal number and possible configurations of limit cycles of planar polynomial vector fields of degree n.}
\subsection{Short English statement}
How many isolated periodic motions can a planar polynomial differential equation have, and how can those cycles be arranged?
% BEGIN RESEARCH ADDITION 40
\subsection{Research attempt}
\paragraph{Current frontier.}
\textit{Literature check: 27 September 2026.}
The distinction between finiteness for one vector field and a bound uniform
in all its degree-$n$ coefficients is essential. The latter is not supplied
by the former. The formulation reference studies local cyclicity through
Melnikov functions and Lyapunov constants, rather than the entire global
configuration problem; its journal record gives a degree-six local lower
bound of 44. \sref{40}{1}, \eref{40}{1}.
The 2024 Buzzi--Novaes critique rules out the advertised quadratic formula
for $H(n)$: the established lower-growth estimate is already of order
$n^2\log n$. Thus a newspaper report about that formula is not evidence
that the present target has been solved. \sref{40}{2}.
A 2026 manuscript investigates separable Chebyshev replication and its
quadratic ceiling when no additional cycles are created. It is a useful
construction-level comparison, not a general upper bound. \eref{40}{2}.

\paragraph{Attack route.}
There are three distinct possibilities: construct increasingly many cycles
at a fixed degree, obtain a uniform zero bound for suitable return maps,
or cover the compactified coefficient space by neighborhoods with uniformly
bounded cyclicity. The first must keep the degree fixed; the second must
handle degenerate return maps; and the third needs local bounds at every
limiting vector field, including those with centers, singular curves, and
behavior at infinity. We pursue a concrete angular-return calculation,
then test whether its uniformity survives degeneration. The computations
below are deductions for this notebook; elementary ingredients and the
restricted conclusions are not claimed to be historically new.

\paragraph{Attempt: an exactly solvable diagnostic family.}
Set $s=x^2+y^2$ and consider
\[
 P=xF(s)-yG(x,y),\qquad Q=yF(s)+xG(x,y),
 \tag{40.1}
\]
where $F$ is a real polynomial and $G$ is polynomial and strictly positive
on the closed annulus under consideration. In polar coordinates,
\[
 \dot r=rF(r^2),\qquad \dot\theta=G(r\cos\theta,r\sin\theta).
 \tag{40.2}
\]
Consequently every positive root $a$ of $F$ whose circle lies in the
annulus gives a periodic orbit $r=\sqrt a$. These are the only periodic
orbits there: away from a root, $r$ is strictly monotone, and a trajectory
cannot cross an equilibrium of the scalar equation for $r$ by uniqueness.
If $F\not\equiv0$, each such circle is isolated, even for a multiple root.
If $F\equiv0$, the annulus is filled by periodic orbits and none is a limit
cycle. Thus this family has exactly as many limit cycles in the annulus
as $F$ has distinct positive roots there.

For a simple root $a$, let
\[
 T_a=\int_0^{2\pi}\frac{d\theta}
 {G(\sqrt a\cos\theta,\sqrt a\sin\theta)}.
\]
The transverse Floquet multiplier is
\[
 \mu_a=\exp\bigl(2aF'(a)T_a\bigr).                 \tag{40.3}
\]
Indeed the linearization of $rF(r^2)$ at $\sqrt a$ is
$2aF'(a)$, constant along the orbit. Equivalently, differentiating
$dr/d\theta=rF(r^2)/G$ at the root gives
$2aF'(a)/G$, because the term differentiating $1/G$ is multiplied by
$F(a)=0$. This also checks that the nonconstant angular speed has not been
silently ignored.

Take $G=1$ and
$F(s)=c\prod_{j=1}^{m}(s-a_j)$, with
$0<a_1<\cdots<a_m$ and $c\ne0$. This is a degree-$2m+1$ field
with precisely $m$ nested hyperbolic limit cycles, and their stability
alternates. The example checks existence, isolation, degree, and relative
position explicitly. It is only a lower bound for the unrestricted class;
in particular it is nowhere near the known large lower bounds cited above.

\paragraph{Attempt: first-order angular return maps.}
Now perturb the linear center by arbitrary degree-at-most-$n$ polynomials:
\[
 \dot x=-y+\varepsilon p(x,y),\qquad
 \dot y=x+\varepsilon q(x,y).
\]
On a fixed compact annulus $0<r_-\le r\le r_+<\infty$, write
\[
 R(r,\theta)=p(r\cos\theta,r\sin\theta)\cos\theta+
 q(r\cos\theta,r\sin\theta)\sin\theta,
\]
\[
 T(r,\theta)=\frac{q(r\cos\theta,r\sin\theta)\cos\theta-
 p(r\cos\theta,r\sin\theta)\sin\theta}{r}.
\]
For sufficiently small $|\varepsilon|$, the angular velocity
$1+\varepsilon T$ stays positive and
\[
 \frac{dr}{d\theta}=\frac{\varepsilon R}{1+\varepsilon T}.
\]
Smooth dependence of solutions of this scalar ODE, uniformly on a slightly
larger compact annulus for $0\le\theta\le2\pi$, gives the $C^1$ expansion
\[
 \Pi_\varepsilon(r)-r=\varepsilon M(r)+O(\varepsilon^2),
 \qquad M(r)=\int_0^{2\pi}R(r,\theta)\,d\theta.       \tag{40.4}
\]
For completeness, the zeroth-order trajectory is constant; replacing the
perturbed trajectory in $R$ by this constant changes its value by
$O(\varepsilon)$ by Gronwall's inequality. Multiplying by the exterior
$\varepsilon$ gives the stated remainder. Differentiating the scalar ODE
in the initial radius proves the same bound for the first derivative.
The constants depend on the annulus and bounded coefficients of $p,q$.

Parity now supplies an exact algebraic simplification. A monomial
$x^iy^j$ in $p$ contributes
$r^{i+j}\int\cos^{i+1}\theta\sin^j\theta\,d\theta$, which vanishes
unless $i$ is odd and $j$ even. The analogous term from $q$ survives only
when $i$ is even and $j$ odd. Hence
\[
 M(r)=rA(r^2),\qquad \deg A\le
 \left\lfloor\frac{n-1}{2}\right\rfloor.             \tag{40.5}
\]
If all positive roots of $A$ in the annulus are simple and none is on its
boundary, the implicit function theorem gives exactly one nearby limit
cycle for every such root and no other cycles confined to the annulus,
for sufficiently small nonzero $\varepsilon$. To justify ``no other'',
remove small disjoint neighborhoods of the roots. On the remaining compact
set $|M|$ has a positive minimum, and (40.4) excludes return-map zeros.
Inside each root neighborhood, $M'$ is bounded away from zero, so the
$C^1$ remainder gives uniqueness as well as existence.
An orbit making several full turns would give a periodic point of the
strictly increasing scalar return map. Such a point must be fixed: if
$\Pi(r)>r$ or $\Pi(r)<r$, successive defined iterates remain strictly
ordered and cannot return to $r$. Thus multiple-turn periodic points do
not evade the zero count.


As a check, choose $p=x(s-1)(s-4)$ and $q=y(s-1)(s-4)$.
Then $M(r)=2\pi r(r^2-1)(r^2-4)$, and the two cycles at radii
1 and 2 are in fact exact for every nonzero $\varepsilon$. Their transverse
exponents have opposite signs. The accompanying symbolic script verifies
the polar identities and these signs exactly.

\paragraph{Stress test: why this does not produce a degree bound.}
First, (40.5) bounds zeros of the \emph{first nonzero displacement term
only when that term is actually the first-order one}. Taking
$p=-yB(s)$ and $q=xB(s)$ gives $R=0$ and $M\equiv0$: the perturbation
changes only angular speed, leaving a center. More generally $M\equiv0$
does not imply that the full return map is the identity, and higher-order
terms can control the bifurcation. Bounding their order of appearance and
their zeros uniformly is not supplied by the calculation. This is precisely
why the local-cyclicity literature uses Lyapunov constants and higher-order
information. \sref{40}{1}, \eref{40}{1}.

Second, the annulus estimates are not uniform as $r_-\downarrow0$,
$r_+\uparrow\infty$, or the angular velocity approaches zero. Cycles may
approach a singular point, a separatrix configuration, or infinity instead
of one of the regular return-map charts. A fixed compact annulus is not a
compactification of all these phenomena.

Here is an exact compactness reduction which isolates the omitted step.
Normalize a nonzero degree-at-most-$n$ coefficient vector to Euclidean norm
one; multiplying a vector field by a positive constant preserves all its
orbits. The normalized vectors form a compact sphere $K_n$. Suppose that
for every $v\in K_n$ there were an open neighborhood $U_v$ and an integer
$B_v$ such that \emph{every} vector field with coefficient vector in $U_v$
had at most $B_v$ limit cycles in the entire plane. A finite subcover gives
\[
 H(n)\le\max_{1\le j\le N}B_{v_j}<\infty.           \tag{40.6}
\]
This implication is proved, but its hypothesis is unproved at degenerate
fields. Pointwise finiteness does not imply this local boundedness: the
function on $\{0\}\cup\{1/j:j\ge1\}$ given by $b(0)=0$ and $b(1/j)=j$
is finite at every point of a compact space and has no finite bound.
No upper semicontinuity of the cycle count has been established here.

Third, even (40.6) would address only finiteness, not the target's request
for the exact maximum and all realizable configurations. Nested circles
in (40.1) do not classify nonnested or mixed configurations. A replication
scheme whose degree increases at every step cannot disprove finiteness at
one fixed degree; compare the explicitly degree-increasing 2026
construction. \eref{40}{2}.

\needspace{14\baselineskip}
\paragraph{Outcome.}
\textbf{Outcome: Exploratory approach with unresolved gap.}
We obtained an exact cycle count and Floquet formula in an angular-radial
subclass, a uniform first-order return-map calculation on regular compact
annuli, and a precise compactness implication for global finiteness.
The stress tests identify loss of control when the first Melnikov function
vanishes or return charts degenerate. There is no new global bound on
$H(n)$, no fixed-degree unbounded construction, and no classification of
all configurations. A resolution would need uniform control at every
limiting field and, beyond finiteness, a separate sharpness and realization
analysis. Historical novelty of the diagnostic lemmas is not asserted.
% END RESEARCH ADDITION 40
\subsection{Sources}
\begin{itemize}\raggedright
\item[\textnormal{[S1]}]\hypertarget{source:40:1}{} Luiz Fernando da Silva Gouveia and Joan Torregrosa, The local cyclicity problem. Melnikov method using Lyapunov constants, Proceedings of the Edinburgh Mathematical Society 65 (2022), 356-375, Introduction, p. 356.
\url{https://recercat.cat/bitstream/handle/2072/532013/MelnikovMethod.pdf?sequence=1}.
{\small\textit{Catalogue formulation source.}}
\item[\textnormal{[S2]}]\hypertarget{source:40:2}{} A note on a recent attempt to solve the second part of Hilbert's 16th Problem.
\url{https://arxiv.org/abs/2411.09594}.
{\small\textit{Catalogue research/status reference; not a proof endorsement.}}
\item[\textnormal{[S3]}]\hypertarget{source:40:3}{} Pesquisadores da Unesp propõem solução para desafio matemático em aberto há mais de um século.
\url{https://jornal.unesp.br/2024/09/24/pesquisadores-da-unesp-propoem-solucao-para-desafio-matematico-em-aberto-ha-mais-de-um-seculo/}.
{\small\textit{Catalogue research/status reference; not a proof endorsement.}}
% BEGIN RESEARCH SOURCES 40
\item[\textnormal{[E1]}]\hypertarget{extra:40:1}{} Luiz F. S. Gouveia and Joan Torregrosa,
\emph{The local cyclicity problem: Melnikov method using Lyapunov constants},
Proceedings of the Edinburgh Mathematical Society 65(2) (2022), 356--375,
DOI 10.1017/S0013091522000128. Author-institution record and abstract,
including the local degree-six bound and the first-order comparison.
\url{https://portalrecerca.uab.cat/en/publications/the-local-cyclicity-problem-melnikov-method-using-lyapunov-consta}.
{\small\textit{Accessible identification of [S1]; consulted 27 September 2026.
The original repository link did not load in this research run.}}
\item[\textnormal{[E2]}]\hypertarget{extra:40:2}{} Olimjon Eshkobilov, Shirali Kadyrov,
and Khudoyor Mamayusupov, \emph{Limit-Cycle Replication via Chebyshev Pullbacks
and a Quadratic Ceiling for Separable Schemes}, arXiv:2604.12883v1,
14 April 2026, Introduction, Theorems 1--2 and Section 5.
\url{https://arxiv.org/html/2604.12883v1}.
{\small\textit{Research manuscript, not treated as an independently verified
resolution; consulted 27 September 2026.}}
% END RESEARCH SOURCES 40
\end{itemize}

\end{document}
